% Options for packages loaded elsewhere
\PassOptionsToPackage{unicode}{hyperref}
\PassOptionsToPackage{hyphens}{url}
\documentclass[
  ignorenonframetext,
]{beamer}
\newif\ifbibliography
\usepackage{pgfpages}
\setbeamertemplate{caption}[numbered]
\setbeamertemplate{caption label separator}{: }
\setbeamercolor{caption name}{fg=normal text.fg}
\beamertemplatenavigationsymbolsempty
% remove section numbering
\setbeamertemplate{part page}{
  \centering
  \begin{beamercolorbox}[sep=16pt,center]{part title}
    \usebeamerfont{part title}\insertpart\par
  \end{beamercolorbox}
}
\setbeamertemplate{section page}{
  \centering
  \begin{beamercolorbox}[sep=12pt,center]{section title}
    \usebeamerfont{section title}\insertsection\par
  \end{beamercolorbox}
}
\setbeamertemplate{subsection page}{
  \centering
  \begin{beamercolorbox}[sep=8pt,center]{subsection title}
    \usebeamerfont{subsection title}\insertsubsection\par
  \end{beamercolorbox}
}
% Prevent slide breaks in the middle of a paragraph
\widowpenalties 1 10000
\raggedbottom
\AtBeginPart{
  \frame{\partpage}
}
\AtBeginSection{
  \ifbibliography
  \else
    \frame{\sectionpage}
  \fi
}
\AtBeginSubsection{
  \frame{\subsectionpage}
}
\usepackage{iftex}
\ifPDFTeX
  \usepackage[T1]{fontenc}
  \usepackage[utf8]{inputenc}
  \usepackage{textcomp} % provide euro and other symbols
\else % if luatex or xetex
  \usepackage{unicode-math} % this also loads fontspec
  \defaultfontfeatures{Scale=MatchLowercase}
  \defaultfontfeatures[\rmfamily]{Ligatures=TeX,Scale=1}
\fi
\usepackage{lmodern}
\ifPDFTeX\else
  % xetex/luatex font selection
\fi
% Use upquote if available, for straight quotes in verbatim environments
\IfFileExists{upquote.sty}{\usepackage{upquote}}{}
\IfFileExists{microtype.sty}{% use microtype if available
  \usepackage[]{microtype}
  \UseMicrotypeSet[protrusion]{basicmath} % disable protrusion for tt fonts
}{}
\makeatletter
\@ifundefined{KOMAClassName}{% if non-KOMA class
  \IfFileExists{parskip.sty}{%
    \usepackage{parskip}
  }{% else
    \setlength{\parindent}{0pt}
    \setlength{\parskip}{6pt plus 2pt minus 1pt}}
}{% if KOMA class
  \KOMAoptions{parskip=half}}
\makeatother
\setlength{\emergencystretch}{3em} % prevent overfull lines
\providecommand{\tightlist}{%
  \setlength{\itemsep}{0pt}\setlength{\parskip}{0pt}}
% Auto-generated by infrastructure/rendering/_pdf_latex_helpers.py::extract_math_font_preamble.
% Minimal header passed to Pandoc via -H so Beamer slides pick up the same math font
% as the combined-PDF path. Adjust the manuscript's preamble.md to change the font globally.
\usepackage{unicode-math}
\setmathfont{latinmodern-math.otf}

% Slide layout defaults for warning-clean scientific decks.
\usepackage{etoolbox}
\IfFileExists{xurl.sty}{\usepackage{xurl}}{}
\IfFileExists{seqsplit.sty}{\usepackage{seqsplit}}{\newcommand{\seqsplit}[1]{#1}}
\protected\def\breakseq#1{\seqsplit{#1}}
\protected\def\breaktt#1{\begingroup\ttfamily\seqsplit{#1}\endgroup}
\setlength{\emergencystretch}{6em}
\tolerance=5000
\hbadness=10000
\hfuzz=1pt
\setlength{\tabcolsep}{2pt}
\AtBeginEnvironment{longtable}{\tiny\renewcommand{\arraystretch}{0.86}\setlength{\tabcolsep}{1pt}}
\AtBeginEnvironment{tabular}{\tiny\renewcommand{\arraystretch}{0.86}\setlength{\tabcolsep}{1pt}}

% Natbib and cross-reference fallbacks — slides don't load natbib
% or cleveref, but combined-PDF manuscript prose may emit these
% commands. The fallback renders citations as a bracketed key list
% and cross-references as detokenized labels so slides don't fail on
% undefined control sequences or raw underscores. \providecommand is
% a no-op if packages load later.
\providecommand{\citep}[1]{[#1]}
\providecommand{\citet}[1]{#1}
\providecommand{\citealp}[1]{#1}
\providecommand{\citeauthor}[1]{#1}
\providecommand{\citeyear}[1]{#1}
\providecommand{\cref}[1]{\texttt{\detokenize{#1}}}
\providecommand{\Cref}[1]{\texttt{\detokenize{#1}}}

\newtheorem{proposition}{Proposition}
\newtheorem{hypothesis}{Hypothesis}
\usepackage{bookmark}
\IfFileExists{xurl.sty}{\usepackage{xurl}}{} % add URL line breaks if available
\urlstyle{same}
% fallback for those not using the hyperref driver hyperxmp:
\makeatletter
\@ifundefined{xmpquote}{\newcommand{\xmpquote}[1]{#1}}{}
\makeatother
\hypersetup{
  hidelinks,
  pdfcreator={LaTeX via pandoc}}

\author{\texorpdfstring{}{}}
\date{}

\begin{document}

\section{Discussion and Conclusion}\label{sec:discussion}

\begin{frame}[allowframebreaks]{Discussion and Conclusion}
DuckRabbit's principal result is a reproducibility boundary rather than
a new psychophysical finding. The same typed request can be regenerated,
hashed, encoded, decoded, and audited, while the evidence graph and
observer layer remain explicit about what the package cannot infer. This
makes the atlas useful for stimulus construction, source comparison, and
preregistration without treating a rendered image or sound as behavioral
evidence.

The scholarly contribution is correspondingly modest but operational.
Rather than flattening visual, auditory, temporal, and multisensory
families into a single ``illusion'' label, the package keeps mechanism,
perceptual signature, requirements, evidence role, engineering basis,
and implementation status separate. A primary demonstration, a review, a
theoretical account, and a source describing the DuckRabbit
implementation answer different questions. The appendix and source-data
sidecars make those distinctions inspectable at the same time as the
generated media.

The work is best understood as a research-software artifact with a
bounded methods contribution. Its novelty claim is not that it discovers
a new illusion or resolves a disputed mechanism; it is that a
heterogeneous stimulus catalog can be represented as typed,
reproducible, source-linked, and release-audited objects. This framing
is consistent with software-citation and FAIR guidance, which treats
versioned software, metadata, provenance, and reuse conditions as part
of the research record \breakseq{[@smith2016software; @wilkinson2016fair;
@lamprecht2020fairsoftware]}.

The release is authored by Daniel Ari Friedman of the Active Inference
Institute, and its intended public home is
\href{https://github.com/docxology/DuckRabbit}{the DuckRabbit public
GitHub repository}. That repository statement is part of the software's
citation identity; until the external handoff occurs, the private
sidecar and its generated bundle are the authoritative release
candidate.
\end{frame}

\begin{frame}[allowframebreaks]{Conclusion}
\protect\phantomsection\label{sec:conclusion}
DuckRabbit v0.5.0 turns multimodal illusion generation into an auditable
typed pipeline. A request produces a canonical artifact, objective
metrics, optional delivery files, decoded inspection, and a versioned
evidence-bounded manifest. The catalog now records not only what is
implemented, but also what the literature supports and what remains
unvalidated.

The main scientific contribution is a boundary: deterministic stimulus
facts are reproducible software outputs, while perceptual effects are
hypotheses that require observer conditions and data. This boundary
permits broad generator coverage without overclaiming. Future additions
should contribute a typed parameter contract, a literature record, an
engineering-fidelity statement, objective invariants, generated
documentation, and tests before entering the implemented registry.
Promotion is a release decision about software readiness, not a verdict
on whether a perceptual phenomenon is real.

The synthetic-psychophysics layer extends this boundary without crossing
it. A transparent feature observer makes model inputs, weights,
calibration, and output hashes inspectable, so the full orchestration
can be tested today. Its small curve is a diagnostic of that specified
model, not evidence that humans share its feature map or response
function. Empirical observer work remains a separate, ethically and
methodologically governed stage.
\end{frame}

\end{document}
